\documentclass[12pt, a4paper]{article}
\usepackage[utf8]{inputenc}
\usepackage{titling}
\usepackage{graphicx}
\graphicspath{ {./images/} }
\usepackage[spanish]{babel}

\usepackage{array}
\usepackage{tabularx}
\usepackage{lipsum}
\usepackage{hyperref}

\usepackage{amsthm}
\newtheorem{theorem}{Teorema}[section]
\newtheorem{lemma}[theorem]{Lema}

\usepackage{minted}
\newminted{rust}{}
\newmintinline{rust}{}

\title{
    \includegraphics[height=0.2\textwidth]{uns-logo}
    \includegraphics[height=0.2\textwidth]{dcic-logo} \\
    UNS Algoritmos y Complejidad 2021 Actividad 14 \\
    \large{Aplicaciones de Componentes Fuertemente Conexos en Compiladores}
    }
\author{Iñaki Garay L.U. 67387}
\date{Julio 2021}

\begin{document}

\maketitle

% á é í ó ú ñ

\section{Introducción}
Este informe documenta tres aplicaciones del algoritmo para encontrar \textbf{Componentes Fuertemente Conexos} en el contexto de compiladores de lenguajes de programación.

Las primeras dos aplicaciones, la planificación de paralelismo a nivel de instrucciones para aprovechar recursos de microprocesador (\textit{instruction-level parallelism scheduling}), y el ánalisis de punteros sensible al contexto, son aplicaciones tradicionales y conocidas. 
El último ejemplo, el uso de CFCs en el grafo de flujo de llamadas (\textit{call-flow graph}) para mejorar el análisis estatico en el compilador de Rust, es más novedoso.

\section{Componentes Fuertemente Conexos}

A graph is said to be strongly connected if every vertex in the graph is reachable from every other vertex.

\section{Planificación de Paralelismo a Nivel de Instrucciones}

\subsection{Contexto}

\subsection{Aplicación}

\section{Análisis de Punteros Sensible al Contexto}

\subsection{Contexto}

\subsection{Aplicación}

\section{El \textit{borrow-checker} de Rust}

\subsection{Contexto}

El lenguaje de programación Rust es un lenguaje moderno, de propósito general, compilado, con tipado estático, que ocupa el mismo espacio que C++. 
Tiene una caracteristica particular: esta diseñado para garantizar que el acceso a memoria sea \textit{seguro}, (\textit{i.e.} no es posible desreferenciar punteros nulos, tener punteros colgantes, o liberar dos veces el mismo puntero) a pesar de no utilizar un recolector de basura o requerir manejo de memoria manual. 

Para esto se vale de una formalizacion del concepto de pertenencia (\textit{ownership}) en su sistema de tipos. 
Que un objeto en memoria pertenezca a otro generalmente implica que el dueño decide cuando se libera la memoria asociada al objeto: al destruirse el dueño, todas sus pertenencias se destruyen tambien. 

% En Rust, este patron admeas se asocia con el de la mutabilidad. 
% La motivacion detras de este concepto es la observacion que la mayoria de los errores relacionados a memoria ocurren cuando valores en memoria ...

% \begin{center}
% \begin{tabular}{ |r|c|c|c| } 
% \hline
% Tipo             & Ownership         & Alias & Mutable \\ 
% \hline
% \texttt{T}       & Owned             &    No & No      \\ 
% \texttt{mut T}   & Owned             &    No & Si      \\
% \texttt{\&T}     & Shared Reference  &    Si & No      \\ 
% \texttt{\&mut T} & Mutable Reference &    No & Si      \\ 
% \hline
% \end{tabular}
% \end{center}

Este sistema hace valer tres reglas en todo momento: 

\begin{enumerate}
\item Cada valor tiene una variable denominada su \textit{dueña}.
\item Cada valor solo puede tener una dueña a la vez.
\item Cuando una dueña sale de alcance, su valor sera destruido (\textit{i.e.} su memoria es liberada y el valor no podrá ser utilizado).
\end{enumerate}

La pertenencia puede ser transferida, o se puede efectuar un "prestamo": ...

Estas reglas aseguran varias propiedades:

\begin{enumerate}
\item Todas las variables estan inicializadas antes de ser usadas. 
\item No se puede mover un mismo valor dos veces.
\item No se puede mover un valor que ha sido prestado.
\item No se puede acceder una locacion de memoria mientras ha sido prestada de manera mutable (salvo a traves de una referencia mutable).
\item No se puede mutar una locacion de memoria mientras este tomada en prestamo de manera inmutable. 
\end{enumerate}

El compilador de Rust convierte el codigo fuente sucesivamente en distintas representaciones intermedias, hasta llegar al MIR (\textit{Mid-level Intermediate Representation}), una representacion intermedia basada en un grafo de flujo de llamadas. 
Sobre esta reprepresentacion en forma de grafo se verifican las propiedades de las cuales es responsable el \textit{borrow-checker}, la fase del compilador que hace valer las reglas de pertenencia.

A modo de ejemplo, ...


Otro ejemplo...

\begin{minted}[fontsize=\small]{rust}
fn main() {
    let mut vec = Vec::new();
    vec.push(1);
    vec.push(2);
}
\end{minted}

\url{https://play.rust-lang.org/?gist=30074856e62e74e91f06abd19bd72ece&version=stable}

\begin{minted}[fontsize=\small]{rust}
fn main() -> () {
    let mut _0: ();                      // return place in scope 0 at src/main.rs:1:11: 1:11
    let mut _1: std::vec::Vec<i32>;      // in scope 0 at src/main.rs:2:9: 2:16
    let _2: ();                          // in scope 0 at src/main.rs:3:5: 3:16
    let mut _3: &mut std::vec::Vec<i32>; // in scope 0 at src/main.rs:3:5: 3:8
    let _4: ();                          // in scope 0 at src/main.rs:4:5: 4:16
    let mut _5: &mut std::vec::Vec<i32>; // in scope 0 at src/main.rs:4:5: 4:8
    scope 1 {
        debug vec => _1;                 // in scope 1 at src/main.rs:2:9: 2:16
    }

    bb0: {
        _1 = Vec::<i32>::new() -> bb1;   // scope 0 at src/main.rs:2:19: 2:29
                                         // mir::Constant
                                         // + span: src/main.rs:2:19: 2:27
                                         // + user_ty: UserType(0)
                                         // + literal: Const { ty: fn() -> std::vec::Vec<i32> {std::vec::Vec::<i32>::new}, val: Value(Scalar(<ZST>)) }
    }

    bb1: {
        _3 = &mut _1;                    // scope 1 at src/main.rs:3:5: 3:8
        _2 = Vec::<i32>::push(move _3, const 1_i32) -> [return: bb2, unwind: bb5]; // scope 1 at src/main.rs:3:5: 3:16
                                         // mir::Constant
                                         // + span: src/main.rs:3:9: 3:13
                                         // + literal: Const { ty: for<'r> fn(&'r mut std::vec::Vec<i32>, i32) {std::vec::Vec::<i32>::push}, val: Value(Scalar(<ZST>)) }
    }

    bb2: {
        _5 = &mut _1;                    // scope 1 at src/main.rs:4:5: 4:8
        _4 = Vec::<i32>::push(move _5, const 2_i32) -> [return: bb3, unwind: bb5]; // scope 1 at src/main.rs:4:5: 4:16
                                         // mir::Constant
                                         // + span: src/main.rs:4:9: 4:13
                                         // + literal: Const { ty: for<'r> fn(&'r mut std::vec::Vec<i32>, i32) {std::vec::Vec::<i32>::push}, val: Value(Scalar(<ZST>)) }
    }

    bb3: {
        _0 = const ();                   // scope 0 at src/main.rs:1:11: 5:2
        drop(_1) -> bb4;                 // scope 0 at src/main.rs:5:1: 5:2
    }

    bb4: {
        return;                          // scope 0 at src/main.rs:5:2: 5:2
    }

    bb5 (cleanup): {
        drop(_1) -> bb6;                 // scope 0 at src/main.rs:5:1: 5:2
    }

    bb6 (cleanup): {
        resume;                          // scope 0 at src/main.rs:1:1: 5:2
    }
}
\end{minted}


\subsection{Aplicación}

\section{Referencias}

\begin{itemize}
\item Compilers. Principles, Techniques and Tools (Ahi, Sethi, Ullman, 2nd Edition, 2009), Chapter 10 Instruction-Level Parallelism.
\item Compilers. Principles, Techniques and Tools (Ahi, Sethi, Ullman, 2nd Edition, 2009), Chapter 10 Instruction-Level Parallelism.
\item The Rust Programming Language, Chapter 4 Understanding Ownership 
      ( \url{https://doc.rust-lang.org/book/ch04-00-understanding-ownership.html} ). 
\item The Rustc Developer Guide, Chapter 44 The Borrow-Checker 
      ( \url{https://rustc-dev-guide.rust-lang.org/borrow_check.html} ). 
\end{itemize}

\end{document}